\section{Evaluation}
\label{sec:evaluation}

\subsection{Experimental setup}
\label{sec:setup}

\textbf{Hardware.} Three graphics adapters across two machines, both running Windows 11 with DirectX 12 and DDR4-3200 memory:

\begin{sloppypar}
\begin{itemize}[leftmargin=1.4em,itemsep=1pt,topsep=2pt]
    \item \textit{Laptop} -- ASUS TUF Gaming F15 FX506HC: Core i5-11400H (Tiger Lake-H, 6C/12T, 4.1 GHz), 16 GB dual channel, 144 Hz panel. Its hybrid configuration exposes \emph{Intel UHD} (\texttt{ANGLE\allowbreak{} (Intel, Intel(R) UHD Graphics (0x00009A68)\dots)}, sharing system memory) and \emph{RTX 3050} (\texttt{ANGLE\allowbreak{} (NVIDIA, \dots RTX 3050 Laptop GPU (0x000025A2)\dots)}, 4 GB dedicated).
    \item \textit{Desktop} -- MSI Z390-A Pro: Core i7-8700K (Coffee Lake-S, 6C/12T, 4.7 GHz), 32 GB, 60 Hz display, \emph{GTX 1080} (\texttt{ANGLE\allowbreak{} (NVIDIA, \dots GTX 1080 (0x00001B80)\dots)}, 8 GB dedicated).
\end{itemize}
\end{sloppypar}

Comparing Intel UHD with the RTX 3050 changes only the GPU within the same host and session. Comparing either with the GTX 1080 changes the host, while the operating system family, memory speed, browser, and engine build remain matched. Each adapter was explicitly assigned per application through the Windows graphics settings, and the hardware was identified throughout by the renderer string recorded in every measurement.

\textbf{Software and build.} The campaign ran in Microsoft Edge 151.0.4129.107 at device pixel ratio 1.0 in a fixed window, with the laptop on AC power. Uncapped runs used a separate clean profile started with \texttt{--disable-gpu-vsync --disable-frame-rate-limit}. Every row carries the engine build stamp, identical throughout (\texttt{engine\_builtAt = 2026-08-24T18:43:54.699Z}). 

\textbf{Benchmark scenes.} Every configuration loads a real scenario archive, the same ZIP the platform serves. Table~\ref{tab:scenes} gives the fingerprint recorded in every row, identical on all three adapters.

\begin{table}[H]
\centering\small\setlength{\tabcolsep}{4pt}
\caption{Benchmark scenes, from the fingerprint columns carried by every measured row.}
\label{tab:scenes}
\begin{tabular}{lrrrr}
\toprule
Scene & GameObjects & Triangles & Draw calls & Textures \\
\midrule
1 -- primitive grid, $N=100$   & 104   & 23{,}194    & 87    & 1 \\
1 -- $N=500$   & 504   & 111{,}926   & 427   & 1 \\
1 -- $N=1000$  & 1{,}004 & 213{,}450 & 812   & 1 \\
1 -- $N=5000$  & 5{,}004 & 1{,}033{,}214 & 3{,}931 & 1 \\
2 -- high-polygon textured model & 14  & 436{,}810   & 4     & 4 \\
3 -- Solar System scenario     & 21    & 37{,}132    & 15    & 15 \\
\bottomrule
\end{tabular}
\end{table}

\textit{Scene 1} is a grid of $N$ primitives over a floor plane sharing meshes and materials, in which every object rotates every frame. Object count is content rather than a flag, so the sweep consists of four archives differing only in that constant; each adds a camera, two lights and a floor, so every Scene 1 axis reports the measured \texttt{objects} column rather than the nominal $N$. \textit{Scene 2} is a single PBR model with base-colour, normal and metallic-roughness maps, all $2048^2$, of which only the base-colour map is de-embedded from the GLB into the engine's texture path; texture-cap results for this scene therefore concern that map alone. \textit{Scene 3} is the deployed Solar System scenario, including an $8192\times4096$ equirectangular skybox.

\textbf{Protocol.} A configuration is a set of URL flags. An unattended runner executes $N$ repetitions of $M$ configurations per session, reloading the page between repetitions, because the flags set engine-wide statics that outlive a scene swap: an incomplete reset would not fail visibly but would silently report the previous configuration's numbers under the current name. Each load discards 120 warm-up frames and measures the next 600, giving 29 configurations per adapter and 87 in total. Scene 3 was repeated ten times and the remaining scenes five times, the first repetition of each being discarded, so 9 of 10 and 4 of 5 repetitions contribute to every reported mean.

\textbf{Metrics and reproducibility.} Mean CPU frame time (\texttt{Application.\allowbreak cpuFrameTime}, busy milliseconds on the main thread per loop iteration); wall-clock mean frame time, reported only where the configuration is not refresh-bound; per-phase timings, of which \texttt{Update} and \texttt{Render} are reported; first-render CPU cost; and engine-estimated texture VRAM derived from each live texture's format and dimensions. MB denotes $2^{20}$ bytes. Harness, scenes, measured build, protocol and raw exports are published at \url{https://github.com/HuGoxDP/WebEngineTS-Benchmarks}; configurations are selected entirely through URL parameters, and every row records its own flags, fingerprint, adapter and build stamp.

\subsection{Choice of frame-time metric}
\label{sec:metric}

On two of the three adapters the wall clock measures the display rather than the engine. On the GTX 1080 at 60 Hz, across all 29 configurations spanning 12 to 1{,}033{,}214 triangles, mean frame time lies between 16.666 and 16.668 ms and the frame rate reads 60.00 in every row. On the RTX 3050 at 144 Hz, 24 of 29 configurations sit within 1.5\% of 6.944 ms, and only Scene 1 at 5{,}004 objects falls genuinely below the cap, at 23.956 ms. On Intel UHD, which never reaches the same panel's refresh rate, no configuration is bound. In total 53 of 87 configurations are refresh-rate-bound, so a comparison built on those cells would show every optimization changing nothing in scenes that never dropped below the cap -- an unmeasured result rather than a null one. Wall-clock time would also mean something different on each adapter, whereas CPU frame time means the same on all three.

Removing VSync does not resolve this. On the integrated adapter it destroys the measurement: uncapped, the engine pushes roughly 2000 fps into an adapter sharing system memory, the queue saturates and the browser stalls for seconds, giving a median frame of 0.50 ms against a maximum of 2709 ms, a standard deviation of 110 ms on a mean of 5.21 ms, and a \texttt{Render} phase rising from 0.61 to 4.83 ms. The CPU is blocked rather than computing, and that condition was not repeated. On the discrete adapters uncapped measurement works but largely measures the timer: Scenes 2 and 3 render in 0.14--0.52 ms with a wall-clock coefficient of variation of 4--22\%. Wall-clock frame time is therefore quoted only on Intel UHD, where nothing is pinned, and for Scene 1 at 5{,}004 objects on the RTX 3050.


\textbf{Variance.} Scene fingerprints did not vary between the repetitions of any
configuration, so the controls stated above held throughout. The worst wall-clock
coefficient of variation is 6.1\% (RTX 3050, Scene 1 at 5{,}004 objects). CPU frame time
exceeds 5\% CV in 22 of 87 configurations -- 5 on Intel UHD, 17 on the RTX 3050 and none
on the GTX 1080 -- and where it does, the cause is scale rather than instability: on the
RTX 3050 the measured quantity in Scenes 2 and 3 is 0.79--0.96 ms, close to the timer's
resolution. Standard deviations are given below wherever a comparison depends on them.

Each strategy is consequently reported against the metric it is able to move rather than as a cumulative stack. The stack was also measured and does not support a combined figure: with all three flags enabled, CPU frame time moves $+4.4$\% on Intel UHD, $-4.5$\% on the RTX 3050 and $-2.1$\% on the GTX 1080 relative to baseline. The same flags make the engine slower on one adapter and faster on two, attributing nothing to any individual strategy.

\subsection{GPU-compressed textures}
\label{sec:res-ktx2}

Texture VRAM is an engine-side estimate derived from format and dimensions and is therefore device-independent. It was verified byte-identical on all three adapters, which is what makes the single table in Table~\ref{tab:ktx2-vram} legitimate.

\begin{table}[H]
\centering\small\setlength{\tabcolsep}{5pt}
\caption{Scene 3 -- estimated texture VRAM, crossed with the texture size cap. Identical on all three adapters; formats are the transcode results recorded in the same rows.}
\label{tab:ktx2-vram}
\begin{tabular}{llrl}
\toprule
Texture cap & Compressed & Texture VRAM (MB) & Formats \\
\midrule
none    & off & 208.0 & \texttt{RGBA8:12} \\
none    & on  & \textbf{69.3} & \texttt{BC7:12} \\
2048    & off & 88.0  & \texttt{RGBA8:12} \\
2048    & on  & 69.3  & \texttt{BC7:12} \\
\bottomrule
\end{tabular}
\end{table}

Uncapped, compression reduces estimated texture VRAM by 66.7\%, from 208.0 to 69.3 MB. This is the only one of the four pairs whose arms differ solely in texture format. The recorded formats establish that the transcode occurred: \texttt{BC7} for all twelve live textures on all three adapters against \texttt{RGBA8} uncompressed, confirmed by a single-cube transcoder check on each adapter. Without that column, an asset silently falling back to an uncompressed upload would be indistinguishable from one on which compression saved less. With the cap at 2048 the reported figure also falls, but that pair is not a memory measurement, for the reason given below.

\begin{table}[H]
\centering\small\setlength{\tabcolsep}{3pt}
\caption{Scene 3 -- what compression costs and saves per adapter. First render is with no texture cap, where the arms differ only in format; CPU frame time is given at both cap settings, VSync on.}
\label{tab:ktx2-effects}
\begin{tabular}{lrll}
\toprule
Adapter & \makecell[r]{First render\\no cap (ms)} & \makecell[l]{CPU frame\\no cap} & \makecell[l]{CPU frame\\cap 2048} \\
\midrule
Intel UHD & 951.80 $\rightarrow$ 412.12 \ ($-56.7$\%) & 0.886 $\rightarrow$ 0.895 & 0.831 $\rightarrow$ \textbf{0.918} (+10.5\%) \\
RTX 3050  & 831.38 $\rightarrow$ 275.72 \ ($-66.8$\%) & 0.793 $\rightarrow$ 0.809 & 0.797 $\rightarrow$ 0.802 (+0.6\%) \\
GTX 1080  & 493.40 $\rightarrow$ 150.55 \ ($\mathbf{-69.5}$\%) & 0.272 $\rightarrow$ 0.274 & 0.264 $\rightarrow$ 0.272 (+3.0\%) \\
\bottomrule
\end{tabular}
\end{table}

Compression removes more than half of the cost of the first rendered frame on every adapter, from $-56.7$\% to $-69.5$\% (Table~\ref{tab:ktx2-effects}), reproduced uncapped at $-67.6$\% and $-69.1$\%. The engine no longer decodes an $8192\times4096$ JPEG on the CPU and uploads it as RGBA8; the transcoded texture reaches the driver in the form it keeps. This is the largest and most consistent effect in the campaign. Where the cap is already set to 2048 the same comparison gives $-47.2$\% and $-50.0$\% on the discrete adapters but $+2.8$\% on Intel UHD, whose first render has a floor near 415 ms that neither flag removes.

Steady-state frame time behaves differently. Uncapped, compression is free on all three adapters within noise; with the cap set, a double-digit cost appears on the integrated adapter alone -- 0.087 ms against a standard deviation of about 0.02 ms, compared with $+0.6$\% and $+3.0$\% on the discrete cards. Block-compressed data cannot be resized through a canvas, so the cap cannot reach a compressed texture: the panorama stays at full resolution and the renderer derives a $6\times4096^2$ uncompressed cube render target instead of a $6\times1024^2$ one. Adapters with 4 and 8 GB of dedicated VRAM absorb that; the integrated adapter takes it from the same dual-channel pool the application is using, and on its wall clock, which is not refresh-bound, the configuration costs $+14.0$\% (14.456 $\rightarrow$ 16.476 ms). The effect is an interaction between two flags rather than a property of compression.

The estimate accounts for live textures but not for the derived cube render target, and the separate render-target estimate covers only shadow maps and post-processing buffers. For Scene 3 the omitted term is 384 MB whenever the source reaches the GPU at full size, which is why the A/B was run on both axes. Uncapped, the term is identical in both arms and cancels (reported skybox delta $-96$ MB, true delta $-96$ MB), so the scene-level figure of 208.0 $\rightarrow$ 69.3 MB is valid. Capped, it does not cancel: the reported $+24$ MB stands against a true $+384$ MB, understating the real cost sixteenfold. That configuration is therefore presented as a flag interaction with its associated frame cost and never as a memory result; for a skybox the effective lever is resolution rather than format.

The cap has no effect on textures embedded in a GLB. On Scene 2, estimated texture VRAM is 48.0 MB with no cap and 48.0 MB at 2048 -- the source maps are exactly $2048^2$ and the downscale triggers strictly above the cap -- and 36.0 MB at 1024, where it reaches only the de-embedded base-colour map. The cap governs what the engine itself decodes; content that leaves its maps inside a GLB opts out.

\thisbottomragged
\subsection{Shader pre-compilation}
\label{sec:res-warmup}

\begin{table}[H]
\centering\small\setlength{\tabcolsep}{5pt}
\caption{Scene 2 -- first-render CPU cost (ms), mean $\pm$ standard deviation.}
\label{tab:warmup}
\begin{tabular}{llrrr}
\toprule
Adapter & Texture cap & Warm-up off & Warm-up on & Change \\
\midrule
Intel UHD & none & 149.58 $\pm$ 5.96  & 147.32 $\pm$ 2.93 & $-1.5$\% \\
Intel UHD & 2048 & 147.05 $\pm$ 3.94  & 149.50 $\pm$ 5.82 & $+1.7$\% \\
Intel UHD & 1024 & 162.63 $\pm$ 9.08  & 161.27 $\pm$ 4.39 & $-0.8$\% \\
\midrule
RTX 3050  & none & 169.90 $\pm$ 7.38  & 156.82 $\pm$ 5.56 & $-7.7$\% \\
RTX 3050  & 2048 & 186.15 $\pm$ 16.45 & 157.03 $\pm$ 4.85 & $-15.6$\% \\
RTX 3050  & 1024 & 172.52 $\pm$ 3.27  & 157.20 $\pm$ 7.07 & $-8.9$\% \\
\midrule
GTX 1080  & none & 102.15 $\pm$ 1.44  & 93.27 $\pm$ 2.27  & $-8.7$\% \\
GTX 1080  & 2048 & 99.25 $\pm$ 3.54   & 94.07 $\pm$ 3.30  & $-5.2$\% \\
GTX 1080  & 1024 & 117.00 $\pm$ 3.12  & 104.30 $\pm$ 3.57 & $-10.9$\% \\
\bottomrule
\end{tabular}
\end{table}

Pre-compilation is effective on both discrete adapters and not on the integrated one. Averaged over the three cap settings, first-render CPU cost falls by 10.9\% on the RTX 3050 and 8.4\% on the GTX 1080, against 0.3\% on Intel UHD (Table~\ref{tab:warmup}). All six discrete pairs improve individually and both results are reproduced uncapped ($-6.3$\% and $-8.6$\%), whereas on the integrated adapter there is no systematic direction and every pair lies within one standard deviation. Pre-compilation recovers time in proportion to how much of the first frame was compilation, and on this integrated part that share is evidently small.

One reading of the table should be resisted. The RTX 3050's warmed configurations cluster within 0.4 ms while the unwarmed span 16.2 ms, but that span is driven largely by a single high-variance configuration, and on the GTX 1080 the corresponding spread closes only from 17.8 to 11.0 ms. The mean is the defensible summary here, not the spread.

Scene 3 shows no effect on any adapter: 413.04 $\rightarrow$ 415.93, 278.22 $\rightarrow$ 276.90 and 155.06 $\rightarrow$ 154.48 ms, all within one standard deviation. Its first render is dominated by asset loading, and a first frame must consist of compilation before pre-compilation can remove it.

\subsection{Dirty-flag transform propagation}
\label{sec:res-dirty}

\begin{table}[H]
\centering\small\setlength{\tabcolsep}{4pt}
\caption{Scene 1 -- mean \texttt{Update} phase time per frame (ms), off / on, VSync on. Rows are the measured \texttt{objects} fingerprint.}
\label{tab:dirty-update}
\begin{tabular}{lrrrrrr}
\toprule
& \multicolumn{2}{c}{Intel UHD} & \multicolumn{2}{c}{RTX 3050} & \multicolumn{2}{c}{GTX 1080} \\
\cmidrule(lr){2-3}\cmidrule(lr){4-5}\cmidrule(lr){6-7}
GameObjects & off / on & ratio & off / on & ratio & off / on & ratio \\
\midrule
104     & 0.149 / 0.056 & 2.7$\times$ & 0.144 / 0.078 & 1.8$\times$ & 0.054 / 0.017 & 3.2$\times$ \\
504     & 0.628 / 0.183 & 3.4$\times$ & 0.498 / 0.190 & 2.6$\times$ & 0.221 / 0.050 & 4.4$\times$ \\
1{,}004 & 1.085 / 0.230 & 4.7$\times$ & 1.016 / 0.351 & 2.9$\times$ & 0.425 / 0.099 & 4.3$\times$ \\
5{,}004 & 3.780 / 1.037 & 3.6$\times$ & 3.501 / 1.071 & 3.3$\times$ & 2.136 / 0.680 & 3.1$\times$ \\
\bottomrule
\end{tabular}
\end{table}

The strategy defers transform synchronization from the property setters into a single pre-render flush, so the phase it can move is \texttt{Update}. That phase becomes 1.8 to 4.7 times faster in all twelve pairings (Table~\ref{tab:dirty-update}). The ratio grows monotonically with object count only on the RTX 3050 (1.8 to 3.3$\times$); on the other two adapters it peaks at intermediate counts, so the speed-up should be read as substantial at every scale rather than as scaling with it. The work does not disappear but moves: the \texttt{Render} phase grows in all twelve pairings, by $+7.1$\% to $+20.9$\%, and at 5{,}004 objects by $+18.3$\%, $+15.2$\% and $+20.9$\%.

\begin{table}[H]
\centering\small\setlength{\tabcolsep}{5pt}
\caption{Scene 1 at 5{,}004 objects -- dirty flags off $\rightarrow$ on. Wall clock omitted where refresh-bound.}
\label{tab:dirty-frame}
\begin{tabular}{llll}
\toprule
Condition & Wall clock (ms) & CPU frame time (ms) & Change \\
\midrule
Intel UHD, VSync on & 26.279 $\rightarrow$ 26.541 & 23.874 $\rightarrow$ 24.509 & $+2.7$\% \\
RTX 3050, VSync on  & 23.956 $\rightarrow$ 24.334 & 23.512 $\rightarrow$ 23.894 & $+1.6$\% \\
GTX 1080, VSync on  & refresh-bound               & 11.470 $\rightarrow$ 12.068 & $+5.2$\% \\
RTX 3050, VSync off & 18.704 $\rightarrow$ 21.156 & 18.335 $\rightarrow$ 20.791 & $\mathbf{+13.4}$\% \\
GTX 1080, VSync off & 12.255 $\rightarrow$ 12.554 & 12.089 $\rightarrow$ 12.391 & $+2.5$\% \\
\bottomrule
\end{tabular}
\end{table}

At frame level every one of the five conditions is a cost, from $+1.6$\% to $+13.4$\% (Table~\ref{tab:dirty-frame}). Scene 1 writes one rotation per object per frame, so every transform is dirty in every frame and nothing is skipped: the coalescer pays the bookkeeping and collects none of the saving, which places the whole sweep below its break-even point. The scene is therefore the least favourable workload for this strategy, whose value is proportional to the fraction of a scene that remains static, and the pair of tables should be read together -- the phase-level speed-up is real, and so is the frame-level cost that accompanies it under this workload.

\subsection{Where the frame budget goes}
\label{sec:res-crossdevice}

\begin{table}[H]
\centering\small\setlength{\tabcolsep}{5pt}
\caption{Scene 1 at 5{,}004 objects and 3{,}931 draw calls -- baseline configuration, VSync on.}
\label{tab:cpubound}
\begin{tabular}{lrrrr}
\toprule
Adapter & Wall clock (ms) & CPU frame (ms) & \texttt{Update} (ms) & \texttt{Render} (ms) \\
\midrule
Intel UHD & 26.279 & 23.874 & 3.780 & 18.599 \\
RTX 3050  & 23.956 & 23.512 & 3.501 & 18.573 \\
GTX 1080  & refresh-bound & 11.470 & 2.136 & 8.860 \\
\bottomrule
\end{tabular}
\end{table}

Changing the GPU does not move this scene; changing the host halves it. The two laptop adapters -- integrated and discrete, same processor, same session -- give \texttt{Render} phase times of 18.599 and 18.573 ms, agreeing to within 0.14\%, and CPU frame times within 1.5\% (Table~\ref{tab:cpubound}). The desktop runs the same scene at 8.860 and 11.470 ms, that is 2.10$\times$ and 2.08$\times$ faster. At 3{,}931 draw calls the bottleneck is therefore draw-call submission on the main thread rather than the GPU, which argues for the engine's instancing and static-batching facilities rather than for GPU-side work. A wall-clock comparison would have shown only that the laptop adapters differ by 1.10$\times$, and nothing at all about the desktop, whose frame time is pinned to its display.

\subsection{Threats to validity}
\label{sec:threats}

\textbf{Coverage.} Three adapters over two machines span integrated and discrete graphics and two processor generations, but mobile devices and low-end Chromebooks remain uncharacterized, and three scenes do not span all workloads. Scene 1 is constructed as the least favourable case for dirty-flag transforms and must not be read as characterizing that strategy on content with a large static majority. A single browser was measured deliberately: it is part of the measured system, and holding it fixed is what makes the cross-machine comparisons valid.

\textbf{The desktop-to-laptop gap exceeds the processors' nominal difference.} Operating system family, memory speed, browser and engine build are matched, so the 2.10$\times$ separation is a processor-and-driver difference rather than an uncontrolled one -- but 4.7 GHz against 4.1 GHz on architectures of comparable per-clock throughput does not account for it. Sustained-load power and thermal limits on the mobile part are the natural explanation, and a 600-frame sample falls exactly in that regime, but neither was instrumented. The GPU-invariance result, measured within one host, is not exposed to this.

\textbf{Condition-dependent effect sizes and unreported metrics.} Because most configurations are refresh-bound, wall-clock frame time is unavailable for the greater part of the campaign; obtaining it more widely would require heavier scenes. The cost of dirty-flag transforms is a cost in all five conditions but ranges from $+1.6$\% to $+13.4$\%, the largest with the frame cap removed, so the direction is reportable and the magnitude is not. Texture memory is an estimate from format and dimensions rather than a driver query -- comparable across adapters and verified identical on all three, but omitting the derived cube render target of Scene 3. Loading time, percentile and maximum frame time, and JavaScript heap were recorded but are not reported, being dominated by the refresh cap on two of three adapters; no claim here rests on them.